\section{推荐阅读}

\begin{readingbox}
\textbf{基础}
\begin{itemize}[nosep]
  \item John Holland, \emph{Emergence}，用于理解反复的局部互动如何巩固为更高层级的模式。
  \item Eugene Izhikevich, \emph{Dynamical Systems in Neuroscience}，用于吸引子、盆地与转变语言。
  \item Stanislas Dehaene, \emph{Consciousness and the Brain}，用于从基于状态的神经科学走向实践解释的一条可读桥梁。
\end{itemize}
\textbf{现代研究}
\begin{itemize}[nosep]
  \item \emph{Metastability demystified: the foundational past, the pragmatic
        present and the promising future}（Nature Reviews Neuroscience, 2024，medium）。
  \item \emph{Associative conditioning in gene regulatory network models
        increases integrative causal emergence}（Communications Biology, 2025，medium）。
  \item \emph{Cortical connectivity, local dynamics and stability correlates of
        global conscious states}（Communications Biology, 2025，作为生物学背景属于 strong，迁移到行为层属于 medium）。
  \item \emph{Falling asleep follows a predictable bifurcation dynamic}
        （Nature Neuroscience, 2025，对睡眠状态转变属于 strong）。
\end{itemize}
\textbf{阅读规则}
\begin{itemize}[nosep]
  \item 用现代来源来约束案例中的机制语言。
  \item 用经典教材维持案例的可教学性，并把它们连接到更持久的概念，而不是只依附最近论文术语。
  \item 把前沿意识物理学方案排除在案例层之外；对于日常行为诊断来说，它们过于推测。
\end{itemize}
\end{readingbox}

\begin{summarybox}
案例库现在做的不只是说明框架。它已经成为研究主干与实践迁移之间的一座应用层桥梁。每个案例都识别主导变量，点明相关的路径依赖或状态转变机制，清楚标记证据车道，并提炼出能够超出轶事本身的修复模式。这才让本章像一个分析库，而不只是若干动机故事的合集。
\end{summarybox}